\documentclass[11pt]{article}

\usepackage[margin=1in]{geometry}
\usepackage{amsmath,amssymb}
\usepackage{booktabs}
\usepackage{xcolor}
\usepackage{hyperref}
\usepackage{xurl}
\usepackage{enumitem}
\usepackage{microtype}

\hypersetup{
  colorlinks=true,
  linkcolor=blue!50!black,
  urlcolor=blue!55!black,
  citecolor=blue!50!black
}

\newcommand{\reportbuilddatetime}{2026-06-05 12:43:24 EDT}


\newcommand{\R}{\mathbb{R}}
\newcommand{\GCV}{\operatorname{GCV}}
\newcommand{\RSS}{\operatorname{RSS}}
\newcommand{\df}{\operatorname{df}}
\newcommand{\tr}{\operatorname{tr}}
\newcommand{\argmin}{\operatorname*{arg\,min}}

\title{Local GCV, Synchronized Local Model Selection, and LPS}
\author{%
  \small
  \begin{tabular}{@{}c@{}}
  Source path: \texttt{\string~/current\_projects/geosmooth/split\_handoffs/}\\
  \texttt{lps\_gcv\_synchronized\_selection\_2026-06-05/}\\
  \texttt{lps\_gcv\_synchronized\_selection\_design.tex}
  \end{tabular}%
}
\date{Build: \reportbuilddatetime}

\begin{document}
\maketitle

\section*{Purpose}

This report records a systematic version of the GCV idea for local polynomial
smoothing (LPS).  The motivating question is:

\begin{quote}
Can local weighted-linear-model diagnostics guide the choice of neighborhood
size, kernel, degree, and chart dimension?
\end{quote}

The immediate application is the comparison between
\[
  \texttt{chart.dim = "auto"}
  \qquad\text{and}\qquad
  \texttt{chart.dim = "local.auto"}.
\]
The broader methodological idea is more general.  Every local chart fit in LPS
is a weighted linear regression problem.  Therefore each local fit has
linear-model diagnostics such as residual sum of squares, leverage, effective
degrees of freedom, and generalized cross-validation (GCV).  These quantities
may help choose local modeling parameters in a data-driven way.

The key conceptual turn is this: local GCV should not only be viewed as an
isolated diagnostic for one chart.  It can also be aggregated, smoothed, or
synchronized across overlapping neighborhoods.  This creates a bridge between
ordinary LPS, local parameter selection, and the synchronization philosophy that
motivated S-LPL-TF / SLPLiFT.

\section{The Local Weighted Polynomial Fit}

Let \(x_1,\ldots,x_n\) be observed covariate points and let
\[
  y_i = f_i + \varepsilon_i
\]
be the observed response at point \(i\).  For an anchor point \(i\), LPS chooses
a local support
\[
  N_i \subset \{1,\ldots,n\},
  \qquad m_i = |N_i|.
\]
The local support is usually determined by a neighborhood size, a metric, and a
kernel bandwidth rule.  After constructing a local chart around the anchor,
each neighbor \(j\in N_i\) has local coordinates
\[
  z_{ij}\in \R^{d_i},
\]
where \(d_i\) is the chart dimension used at anchor \(i\).  For polynomial
degree \(p_i\), form the local polynomial design matrix
\[
  X_i \in \R^{m_i\times q_i},
\]
where each row contains the polynomial basis evaluated at \(z_{ij}\).  For
example, with degree one,
\[
  \phi(z) = (1,z_1,\ldots,z_{d_i}),
\]
so \(q_i=d_i+1\).  With degree two, the design also includes all squared and
cross-product terms.

Let
\[
  W_i = \operatorname{diag}(w_{ij}: j\in N_i)
\]
be the diagonal matrix of nonnegative kernel weights, and let
\[
  y_{N_i} = (y_j:j\in N_i)
\]
be the local response vector.  The local weighted least-squares estimate is
\[
  \widehat\beta_i
  =
  \argmin_{\beta\in\R^{q_i}}
  \|W_i^{1/2}(y_{N_i}-X_i\beta)\|_2^2.
\]
When the weighted design has full rank, this is
\[
  \widehat\beta_i
  =
  (X_i^\top W_i X_i)^{-1}X_i^\top W_i y_{N_i}.
\]
In implementation, this should be computed by a stable QR or SVD-based solve
rather than by forming an explicit inverse.

\section{Local Hat Matrix, Degrees of Freedom, and GCV}

For fixed local choices, the fitted values on the local support are a linear
function of the local response vector:
\[
  \widehat y_{N_i}
  =
  H_i y_{N_i},
\]
where
\[
  H_i
  =
  X_i(X_i^\top W_iX_i)^{-1}X_i^\top W_i.
\]
The matrix \(H_i\) is the local weighted hat matrix.  It describes how the
observed local responses are transformed into local fitted values.  Its trace
is the effective local degrees of freedom:
\[
  \df_i = \tr(H_i).
\]
For ordinary unweighted least squares with a full-rank design, \(\df_i\) is the
number of fitted coefficients.  With weights, numerical rank deficiency, or
regularization, \(\df_i\) is the effective model complexity.

The local residual sum of squares is
\[
  \RSS_i
  =
  \|W_i^{1/2}(y_{N_i}-\widehat y_{N_i})\|_2^2.
\]
A local generalized cross-validation score can then be written as
\[
  \GCV_i
  =
  \frac{\RSS_i/m_i}{(1-\df_i/m_i)^2}.
\]
This score penalizes models that fit the local data well only by using too much
local flexibility.  If \(\df_i/m_i\) is close to one, the denominator becomes
small and the score becomes large.  That is exactly the desired behavior when a
local fit is close to being underdetermined or interpolating.

\section{The Caveat and the Opportunity}

The caveat is that \(\GCV_i\) estimates prediction behavior inside the local
neighborhood \(N_i\).  It is not automatically the same thing as the global
held-out prediction error used by dataset-level cross-validation.  Neighboring
local supports overlap heavily, so the local GCV scores
\[
  \GCV_1,\ldots,\GCV_n
\]
are dependent rather than independent observations.

This caveat is also an opportunity.  The dependence is not merely a nuisance;
it is information about the geometry of the smoother.  If nearby anchors share
many data points and produce related local models, then their GCV scores and
their selected local parameters should not be treated as unrelated.  Instead,
we can use overlap and graph adjacency to stabilize local model selection.

This is closely aligned with the motivation behind S-LPL-TF / SLPLiFT.  In
S-LPL-TF, the synchronization term was introduced because neighboring local
models should agree when they make predictions on overlapping regions.  The GCV
view gives a complementary way to express a similar desire: local models should
not only fit well locally, but the field of local choices should be stable and
coherent across the data geometry.

\section{Formulation 1: Global Discrete GCV Selector}

The simplest formulation keeps all tuning parameters global.  Let
\[
  \theta = (k,\kappa,p,d)
\]
denote a candidate parameter setting, where:

\begin{itemize}[leftmargin=*]
  \item \(k\) is the neighborhood/support size;
  \item \(\kappa\) is the kernel type, such as Gaussian or tricube;
  \item \(p\) is the polynomial degree;
  \item \(d\) is the local chart dimension.
\end{itemize}

For a fixed global \(\theta\), every anchor uses the same rule.  Each anchor
still has its own local support, design matrix, and weights, but the model
configuration is global.  Compute
\[
  \GCV_i(\theta),
  \qquad i=1,\ldots,n.
\]
Then select
\[
  \widehat\theta_{\mathrm{global}}
  =
  \argmin_{\theta\in\Theta}
  \sum_{i=1}^{n}\GCV_i(\theta).
\]

This is the most auditable first experiment.  It is a direct analogue of
cross-validation over a candidate grid, but it replaces fold refits with
local linear-model diagnostics.  It should be much cheaper than full K-fold CV
and easier to inspect pointwise.

\paragraph{Robust aggregation.}
Instead of summing raw GCV scores, one may use a robust transformation:
\[
  \widehat\theta_{\mathrm{robust}}
  =
  \argmin_{\theta\in\Theta}
  \sum_{i=1}^{n}\rho\{\GCV_i(\theta)\}.
\]
Here \(\rho\) is increasing but reduces the influence of extreme local failures.
Examples include a winsorized score, a Huberized score, or a trimmed mean.  This
may matter on real geometries where a few local neighborhoods are numerically
unstable.

\section{Formulation 2: Local GCV with a Graph-Smoothed Parameter Field}

The next formulation lets tuning parameters vary by anchor:
\[
  \theta_i = (k_i,\kappa_i,p_i,d_i).
\]
If each anchor independently minimizes its own GCV,
\[
  \widehat\theta_i^{\mathrm{greedy}}
  =
  \argmin_{\theta_i\in\Theta_i}\GCV_i(\theta_i),
\]
then the resulting parameter field may be noisy.  Nearby anchors could choose
very different dimensions or supports because of small local fluctuations.

To borrow strength across nearby anchors, put a graph on the data points:
\[
  G=(V,E),
  \qquad V=\{1,\ldots,n\}.
\]
Then choose the full parameter field
\[
  \theta_{1:n}=(\theta_1,\ldots,\theta_n)
\]
by minimizing
\[
  \sum_{i=1}^{n}\GCV_i(\theta_i)
  +
  \eta
  \sum_{(i,j)\in E}
  D(\theta_i,\theta_j).
\]
Here \(\eta\ge 0\) controls how strongly neighboring local choices are
synchronized, and \(D\) measures disagreement between local parameter choices.

For example, if the most important local choice is chart dimension, one could
use
\[
  D(\theta_i,\theta_j)
  =
  (d_i-d_j)^2.
\]
If support size matters too, one might use
\[
  D(\theta_i,\theta_j)
  =
  \alpha_d(d_i-d_j)^2
  +
  \alpha_k\{\log(k_i)-\log(k_j)\}^2
  +
  \alpha_{\kappa}\mathbf{1}\{\kappa_i\ne \kappa_j\}.
\]
The exact choice of \(D\) is a modeling decision.  The important point is that
this formulation selects a spatially coherent field of local model choices,
rather than a collection of independent local winners.

\section{Formulation 3: Prediction-Synchronized LPS}

The most S-LPL-TF-like formulation synchronizes predictions rather than only
the discrete parameter choices.  Let
\[
  \widehat f^{(i)}(x)
\]
denote the local polynomial model fitted at anchor \(i\), evaluated at a point
\(x\) where that local model is meaningful.  If anchors \(i\) and \(j\) have
overlapping supports, then both local models can often predict values on the
same points.  A prediction-synchronization penalty can be written schematically
as
\[
  \sum_{(i,j)\in\mathcal O}
  \sum_{x_\ell\in N_i\cap N_j}
  \omega_{ij\ell}
  \left\{
    \widehat f^{(i)}(x_\ell)
    -
    \widehat f^{(j)}(x_\ell)
  \right\}^2,
\]
where \(\mathcal O\) is a set of overlapping anchor pairs and
\(\omega_{ij\ell}\) are overlap weights.

A prediction-synchronized LPS criterion could therefore be
\[
  \sum_{i=1}^{n}\GCV_i(\theta_i)
  +
  \eta
  \sum_{(i,j)\in\mathcal O}
  \sum_{x_\ell\in N_i\cap N_j}
  \omega_{ij\ell}
  \left\{
    \widehat f^{(i)}_{\theta_i}(x_\ell)
    -
    \widehat f^{(j)}_{\theta_j}(x_\ell)
  \right\}^2.
\]

This has the same spirit as S-LPL-TF but removes the \(\ell_1\) lifting trend
filtering penalty.  In words: choose local models that have good local
predictive diagnostics and also agree with nearby local models on overlapping
regions.

\section{Relation to S-LPL-TF / SLPLiFT}

S-LPL-TF introduced synchronization because local polynomial lifting models
should not behave as unrelated fits.  Each local model gives predictions on a
support.  When supports overlap, disagreement between those predictions is
evidence that the local models are not forming a coherent global smoother.

The GCV idea complements this.  It says:

\[
  \text{First, ask whether each local model is statistically reasonable.}
\]

Then synchronization says:

\[
  \text{Second, ask whether neighboring local models are mutually coherent.}
\]

So there are two kinds of regularity:

\begin{enumerate}[leftmargin=*]
  \item \textbf{Internal local regularity:} a local model should not use too
  much flexibility for its neighborhood.  GCV, degrees of freedom, leverage, and
  condition diagnostics address this.
  \item \textbf{External overlap regularity:} neighboring local models should
  not make incompatible predictions on shared data.  Synchronization penalties
  address this.
\end{enumerate}

This suggests a larger design principle: local smoothers should balance local
fit, local complexity, and overlap coherence.

\section{Relation to Batched or Randomized Optimization}

The analogy with minibatch optimization in neural networks is useful, but it
should be stated carefully.  The similarity is not that LPS-GCV selection
requires gradient descent.  The similarity is that a large global objective is
assembled from many local losses:
\[
  \mathcal L(\theta)
  =
  \sum_{i=1}^{n}\ell_i(\theta).
\]
In neural-network training, the local losses are usually sample losses, and
minibatches provide noisy but computationally manageable estimates of the full
objective.  In LPS-GCV selection, the local losses are neighborhood-level GCV
scores.  One could imagine estimating
\[
  \sum_{i=1}^{n}\GCV_i(\theta)
\]
using random batches of anchor points rather than all anchors at every
candidate evaluation.

This gives a possible computational strategy:

\begin{enumerate}[leftmargin=*]
  \item sample a batch of anchors \(B\subset\{1,\ldots,n\}\);
  \item evaluate \(\sum_{i\in B}\GCV_i(\theta)\) for candidate \(\theta\);
  \item refine or eliminate candidate parameters based on repeated batches;
  \item only run full-anchor evaluation for the most promising candidates.
\end{enumerate}

The statistical motivation is also similar: do not trust one highly local
quantity too much.  Use many local pieces of evidence, possibly batched,
smoothed, or synchronized, to make more stable choices.

\section{Recommended First Test}

Before building the more ambitious synchronized versions, the first test should
be simple and auditable.

\subsection*{Test 1: Global discrete GCV selector on frozen local-auto batch}

For each frozen dataset and each candidate
\[
  \theta=(k,\kappa,p,d),
\]
compute
\[
  S_{\mathrm{GCV}}(\theta)
  =
  \sum_{i=1}^{n}\GCV_i(\theta).
\]
Then compare the selected candidate
\[
  \widehat\theta_{\mathrm{GCV}}
  =
  \argmin_{\theta} S_{\mathrm{GCV}}(\theta)
\]
to the candidates selected by the existing CV-based auto and local-auto runs.

The output should include:

\begin{itemize}[leftmargin=*]
  \item the GCV-selected support size, kernel, degree, and chart dimension;
  \item the CV-selected support size, kernel, degree, and chart dimension;
  \item Truth RMSE for both selected fits when truth is available;
  \item local degrees-of-freedom ratios \(\df_i/m_i\);
  \item local condition diagnostics;
  \item a comparison between pointwise Truth-RMSE contributions \(c_i\) and
  local GCV diagnostics.
\end{itemize}

The key question is whether red \(c_i\) regions correspond to high GCV, high
leverage, high \(\df_i/m_i\), or unstable local designs.
If the answer is yes, local GCV is not only a selector candidate; it is also an
explanatory diagnostic for why one chart-dimension rule wins or loses.

\section{Possible Next Tests After the First One}

If the global GCV selector is promising, the next natural steps are:

\begin{enumerate}[leftmargin=*]
  \item replace the raw sum of GCV scores by robust summaries, such as trimmed
  means or Huberized sums;
  \item allow chart dimension to vary by anchor while keeping support size and
  kernel global;
  \item add graph-smoothed penalties on the local dimension field;
  \item compare prediction-synchronized LPS to S-LPL-TF without the
  \(\ell_1\) lifting term;
  \item test batched anchor subsets as a scalable approximation for large
  16S-style datasets.
\end{enumerate}

\section*{Conclusion}

Local GCV provides a principled way to use the fact that each local LPS chart is
a weighted linear model.  The simplest version selects one global parameter
setting by minimizing the sum of local GCV scores.  The more interesting
versions treat local GCV scores as a field over the data geometry and stabilize
that field through graph smoothing or prediction synchronization.  This is
conceptually close to the synchronization motivation behind S-LPL-TF, but it
focuses first on local statistical adequacy rather than on trend-filtering
structure.

The recommended first implementation is therefore deliberately modest: compute
local GCV diagnostics on the frozen non-manifold batch and test whether a global
GCV selector agrees with, improves on, or explains the existing CV-selected
auto and local-auto fits.

\appendix
\section{Weighted Hat Matrix Details}

The weighted least-squares problem at anchor \(i\) is
\[
  \min_{\beta}
  (y_{N_i}-X_i\beta)^\top W_i(y_{N_i}-X_i\beta).
\]
Assuming \(X_i^\top W_iX_i\) is nonsingular, the solution is
\[
  \widehat\beta_i
  =
  (X_i^\top W_iX_i)^{-1}X_i^\top W_i y_{N_i}.
\]
The local fitted vector is
\[
  \widehat y_{N_i}
  =
  X_i\widehat\beta_i
  =
  X_i(X_i^\top W_iX_i)^{-1}X_i^\top W_i y_{N_i}.
\]
Thus
\[
  H_i
  =
  X_i(X_i^\top W_iX_i)^{-1}X_i^\top W_i.
\]
The local degrees of freedom are
\[
  \df_i = \tr(H_i).
\]

For numerical implementation, it is often better to work with
\[
  \widetilde X_i = W_i^{1/2}X_i,
  \qquad
  \widetilde y_i = W_i^{1/2}y_{N_i}.
\]
Then the fitted values in the weighted space are obtained by projecting
\(\widetilde y_i\) onto the column space of \(\widetilde X_i\).  If
\(\widetilde X_i\) has singular value decomposition
\[
  \widetilde X_i = U_iD_iV_i^\top,
\]
then the effective rank, leverage diagnostics, and degrees of freedom can be
computed stably from this decomposition.  This is also where rank tolerance and
conditioning diagnostics enter naturally.

\section{Why Ordinary GCV Can Fail Locally}

Local GCV can fail or become unstable when:

\begin{itemize}[leftmargin=*]
  \item \(m_i\) is too small relative to the number of local polynomial
  coefficients;
  \item the local design is nearly rank deficient;
  \item kernel weights concentrate on too few effective observations;
  \item the local neighborhood crosses a singularity, boundary, or mixed
  dimension transition;
  \item the local chart dimension is overestimated.
\end{itemize}

These are not reasons to discard GCV.  They are reasons to record diagnostics
alongside it.  In fact, if high GCV or high \(\df_i/m_i\) aligns with red
pointwise Truth-RMSE contributions, then GCV is doing exactly what we want: it
is identifying local modeling choices that are too flexible or unstable.

\end{document}
